\section{Appendix: Latex Installation}
\label{sec:appendix_latex}  

To be able to use the report generation (View points as well as Evaluation reports), a Latex installation has to be installed on the used workstation. \\

COMPASS starts \texttt{pdflatex} from the \texttt{PATH}. It checks for \texttt{pdflatex} once at startup, so COMPASS has to be started again after the Latex installation. \\

\subsection{Linux}
\label{sec:appendix_latex_linux}

Depending on what Linux operating system is used, the commands will differ. In this section installation instructions for a few selected OSs is provided.

\subsubsection {Ubuntu \& Debian Variants}

\begin{lstlisting}
sudo apt-get install texlive-full 
\end{lstlisting}

\subsubsection {CentOS, Fedora \& AlmaLinux Variants}

\begin{lstlisting}
# AlmaLinux 9 / RHEL 9: enable CRB repository first
sudo dnf config-manager --set-enabled crb
sudo dnf install -y epel-release

sudo dnf install 'texlive-*'
\end{lstlisting}

On AlmaLinux 9 (and similar RHEL 9 derivatives), some LaTeX packages required (e.g. \texttt{ltablex}, \texttt{silence}) are not available in the distribution repositories. These must be installed manually from CTAN into the local texmf tree:

\begin{lstlisting}
# Install ltablex
sudo mkdir -p /usr/share/texlive/texmf-local/tex/latex/ltablex
curl -sL https://mirrors.ctan.org/macros/latex/contrib/ltablex/ltablex.sty \
  | sudo tee /usr/share/texlive/texmf-local/tex/latex/ltablex/ltablex.sty \
  > /dev/null

# Install silence
sudo mkdir -p /usr/share/texlive/texmf-local/tex/latex/silence
curl -sL https://mirrors.ctan.org/macros/latex/contrib/silence/silence.sty \
  | sudo tee /usr/share/texlive/texmf-local/tex/latex/silence/silence.sty \
  > /dev/null

# Update TeX database
sudo texhash /usr/share/texlive/texmf-local
\end{lstlisting}

Verify with \texttt{kpsewhich ltablex.sty} and \texttt{kpsewhich silence.sty}.

\subsection{Windows}
\label{sec:appendix_latex_windows}

On Windows, the Latex distribution \textbf{TeX Live} is used, the same distribution as on Linux:
\begin{itemize}
    \item Download the installer \texttt{install-tl-windows.exe} from \href{https://tug.org/texlive/windows.html}{tug.org/texlive/windows.html}
    \item Run the installer and keep the default installation scheme \textbf{full}
    \begin{itemize}
        \item The full scheme contains all Latex packages the reports use (about 8 GB), no packages have to be installed later
        \item The installer adds the TeX Live \texttt{bin} folder, which contains \texttt{pdflatex.exe}, to the \texttt{PATH}
    \end{itemize}
    \item For workstations without an internet connection, TeX Live is also available as an ISO image with all packages, see \href{https://tug.org/texlive/acquire-iso.html}{tug.org/texlive/acquire-iso.html}
\end{itemize}

The Latex distribution \textbf{MiKTeX} can also be used. Its basic installation does not contain all packages the reports use, they have to be installed in the \textbf{MiKTeX Console} (for example \texttt{adjustbox}, \texttt{fancyhdr}, \texttt{float}, \texttt{geometry}, \texttt{ltablex}, \texttt{pdflscape}, \texttt{silence}). \\

\subsection{Testing}

After installation, the following command should run and generate a similar output (on Windows, the output differs in the version lines):

\begin{lstlisting}
$ pdflatex --version
pdfTeX 3.14159265-2.6-1.40.21 (TeX Live 2020/Debian)
kpathsea version 6.3.2
Copyright 2020 Han The Thanh (pdfTeX) et al.
There is NO warranty.  Redistribution of this software is
covered by the terms of both the pdfTeX copyright and
the Lesser GNU General Public License.
For more information about these matters, see the file
named COPYING and the pdfTeX source.
Primary author of pdfTeX: Han The Thanh (pdfTeX) et al.
Compiled with libpng 1.6.37; using libpng 1.6.37
Compiled with zlib 1.2.11; using zlib 1.2.11
Compiled with xpdf version 4.02
\end{lstlisting}

